\subsection{Checks of the larger-chain numerical method}
\label{sec:scaling-checks}
The Fourier construction in Eq.~\eqref{eq:fourier-covariance} was
compared with the original dense Hamiltonian eigensolve at
$N=8,16,32,128,256$ for all three states. The anchored transition was
compared with the original full complex covariance and singular-value
overlap implementation at $\ell=4$ for every pair. The largest covariance
entry difference was $2.220\!\times\!10^{-15}$ and the largest transition
entry difference was $3.747\!\times\!10^{-15}$. Energies and Fourier applications
to independent test columns were checked as well; all 30 records passed.
At $N=256$, $\ell=4,10$, the new entropy remainder and original
trace-log/Gram formulas agreed within $1.400\!\times\!10^{-12}$ across both CFTs.

At the largest Ising size, Fourier applications additionally verify
$\Gamma^2v=-v$ and antisymmetry on four independent test columns for
each primary; local $XX$ and $Z$ expectations reproduce the analytical
energy per site. These are included in the independent-check table.
Tables~\ref{tab:scaling-numerical-checks} and
\ref{tab:scaling-independent-checks} give the full new production and
precision checks. The support-cutoff criterion was relative variation
below one percent, and the absolute comparison tolerance was
$5\times10^{-12}$. For observables below that absolute scale, independent
high-precision contractions and the noncommuting tiny-Gram example are
essential: a cancellation-damaged zero could otherwise pass an absolute
test. These checks assess numerical accuracy of the finite-chain
construction, not the separate continuum extrapolation. A reported zero
entry difference means identical rounded double-precision outputs;
it does not assert zero error in exact arithmetic or 120-digit accuracy
of the final returned matrices.
\begin{longtable}{lrr}
\caption{Production checks at every added size in Eq.~\eqref{eq:extended-grid}. Relative cutoff sensitivity is $\max_{\tau\in\{10^{-14},10^{-16}\}}|q_\tau-q_{10^{-15}}|/q_{10^{-15}}$, shown here for retained endpoint lengths ($\ell=6,10$ for the dense scan and $\ell=10$ for the earlier extension); the $\ell=4$ checks remain in the raw archive. Trace-log and original full-Gram differences are absolute; they cannot alone certify relative accuracy of extremely small observables. LU residual is $\|KZ-(I-i\Gamma_a)E_A\|_F/\|(I-i\Gamma_a)E_A\|_F$. Positivity errors are $\max(0,-\lambda_{\min})$. All stored tolerances passed.}
\label{tab:scaling-numerical-checks}\\\toprule
Check & Count & Largest error\\\midrule
\endfirsthead
\toprule
Check & Count & Largest error\\\midrule
\endhead
\bottomrule\endfoot
trace & 6390 & $6.661\!\times\!10^{-16}$\\
hermiticity & 6390 & $1.791\!\times\!10^{-77}$\\
positivity & 6390 & $3.324\!\times\!10^{-16}$\\
small\_cutoff\_relative & 4824 & $8.635\!\times\!10^{-9}$\\
small\_trace\_log\_absolute & 4824 & $2.615\!\times\!10^{-13}$\\
local\_LU\_residual & 945 & $1.388\!\times\!10^{-13}$\\
Gram\_positivity & 6390 & $1.240\!\times\!10^{-16}$\\
large\_cutoff\_relative & 4824 & $6.852\!\times\!10^{-9}$\\
large\_full\_Gram\_absolute & 4824 & $1.638\!\times\!10^{-12}$\\
\end{longtable}

\begin{longtable}{lrr}
\caption{Independent numerical checks. Tiny noncommuting Gram errors are relative to a 100-digit matrix-logarithm reference, with cross amplitudes from $10^{-2}$ through $10^{-30}$ (deficits through order $10^{-60}$). Boson diagonal errors are absolute entry differences; transition errors divide the maximum sampled entry difference by the largest reference entry. Sampled entries use 160-digit direct pivoted Pfaffians rather than the 120-digit cached-minor production calculation.}
\label{tab:scaling-independent-checks}\\\toprule
Check & Count & Largest error\\\midrule
\endfirsthead
\toprule
Check & Count & Largest error\\\midrule
\endhead
\bottomrule\endfoot
Ising\_Fourier\_purity & 3 & $5.207\!\times\!10^{-16}$\\
Ising\_Fourier\_skew & 3 & $1.885\!\times\!10^{-18}$\\
Ising\_local\_energy\_per\_site & 3 & $1.110\!\times\!10^{-15}$\\
boson\_diagonal\_high\_precision & 6 & $3.499\!\times\!10^{-78}$\\
boson\_transition\_high\_precision\_relative & 6 & 0\\
tiny\_noncommuting\_Gram & 24 & $8.882\!\times\!10^{-16}$\\
\end{longtable}

